\documentclass{article}
\usepackage{geometry}
\usepackage{graphicx} % Required for inserting images
\usepackage{amsmath, amsthm, amssymb}
\usepackage{parskip}
\newgeometry{vmargin={15mm}, hmargin={24mm,34mm}}
\theoremstyle{definition} 
\newtheorem{definition}{Definition}

\newtheorem{theorem}{Theorem}[section]
\newtheorem{lemma}[theorem]{Lemma}
\newtheorem{proposition}[theorem]{Proposition}
\newtheorem{corollary}{Corollary}[theorem]


\title{MATH210B Homework 7}
\date{March 2025}
\author{Boran Erol}

\begin{document}

\maketitle

\section{Problem 1}

Let $ K/F $ be a field extension and let $\alpha, \beta \in K$. Assume $\alpha$ and $\beta$
are algebraic over $F$, of respective degrees $m$ and $n$.

\begin{lemma}
    Let $m^{\prime}$ be the degree of $\alpha$ over $F(\beta)$. Then, $\beta$ has degree
    $\frac{m^{\prime}n}{m}$ over $F(\alpha)$.
\end{lemma}
\begin{proof}
    Notice that $[F(\alpha, \beta): F(\alpha)][F(\alpha): F] = [F(\alpha, \beta): F(\beta)][F(\beta): F]$.

    Since $[F(\alpha): F] = m$ and $[F(\beta): F] = n$ and $[F(\alpha, \beta): F(\beta)] = m^{\prime}$, the result
    immediately follows. The only lemma we need to continuously apply is that the degree of an element over a field
    is also the degree of the extension.
\end{proof}

\begin{lemma}
    If $m$ and $n$ are coprime, $[F(\alpha, \beta): F] = mn$.
\end{lemma}
\begin{proof}
    By the Tower Law, $[F(\alpha, \beta): F] \leq [F(\alpha): F][F(\beta): F] = mn$.

    Also notice that $m \mid [F(\alpha, \beta): F]$ and $n \mid [F(\alpha, \beta): F]$.
    Since $m$ and $n$ are coprime, $mn \leq [F(\alpha, \beta): F]$.

    We thus conclude the proof.
\end{proof}


\section{Problem 4}

$ y = \frac{x^{2}}{x - 1} $ if and only if $ y(x-1) = x^{2} = 0 $ if and only if $ x^{2} - yx + y = 0 $,
so $ x $ satisfies a quadratic polynomial over $ F(y) $. Notice that $ x $ is not in $ F(y) $
since $ y $ is irreducible in $ F(y) $ and $ y = \frac{x^{2}}{x - 1} $ would contradict this.
Thus, the degree of the extension is $ 2 $.



\section{Problem 5}


Let $ f(x) = x^{3} - 2 \in \mathbb{Q}[x] $. 

Notice that $ f $ is irreducible since we can apply Eisenstein's with $ p = 2 $.
Notice that $ 2 $ doesn't divide the monic term but divides the constant term.
Moreover, $ 4 $ doesn't divide the constant term.

Let $ w = e^{\frac{2\pi i}{3}}  = \frac{-1}{2} + i \frac{\sqrt{3}}{2} $.

Notice that the roots of $ f $ are $ \sqrt[3]{2}, \sqrt[3]{2} w, \sqrt[3]{2} w^{2} $.

Thus, if we add all roots of $ f $, we're also adding $ \sqrt{3} $.

Notice that $ [\mathbb{Q}[\sqrt{3}] : \mathbb{Q}] = 2 $ and $ [\mathbb{Q}\sqrt[3]{2} : \mathbb{Q}] = 3 $.

By Problem 1 on this homework, since $ 2 $ and $ 3 $ are relatively prime, we have that the degree
of the extension is $ 6 $.

\section{Problem 6}

\section{Problem 7}

We present an example of a non-trivial field extension $ K/F $ with isomorphic fields.

Consider $ F = \mathbb{C}[x_{1}, \ldots] $.

Now, consider $ F(y) $. This is a non-trivial field extension since it's a transcendental field extension of $ F $.
However, both of these fields are isomorphic since we can map $ x_{2} \mapsto x_{1}, x_{3} \mapsto x_{2} \cdots $
and map $ x_{1} \mapsto y $.

However, finding a finite field extension would be better. Notice that the extension in Problem 4 satisifies this!



\section{Problem 8}

\begin{lemma}
    Let $ K/F $ be a field extension and $ \alpha \in K $.
    If $ f(\alpha) $ is algebraic over $ F $ for some non-constant $ f \in F[x] $,
    then $ \alpha $ is algebraic over $ F $.
\end{lemma}
\begin{proof}
    If $ f(\alpha) $ is algebraic over $ F $, there exists some $ g \in F[x] $ such that
    $ g(f(\alpha)) = 0 $.  But then $ \alpha $ is a root of $ g \circ f $ and therefore
    algebraic over $ F $.
\end{proof}

\section{Problem 9}

Notice that $ \alpha $ and $ \beta $ are the roots of $ x^{2} - (\alpha + \beta)(x) + \alpha \beta $
which is in $ F(\alpha + \beta, \alpha \beta) $.

\section{Problem 10}

We present an infinite algebraic extension. Consider adding adjoining
all $ p $th roots of unity to $ \mathbb{Q} $ for every prime $ p $. This
is an algebraic extension since every element satisfies $ f(x) = x^{p} - 1 $
for some $ p > 0 $ prime. Moreover, for a $ p $th root of unity $ \zeta $,
we have that $ [\mathbb{Q}(\zeta) : \mathbb{Q}] = p - 1 $. Thus, the field extension is infinite.


\end{document}
